% GENERATED FILE. Edit canonical lessons, then run npm run generate:books.
% canonical-source-hash: fnv1a64:990102c947481530
% canonical-lessons: FR-C277-etape, FR-C277-phase, FR-C277-chronologie, FR-C277-instant, FR-C277-a-l-issue-de

\chapter{Étape par étape}
\label{ch:fr-etape-par-etape}
\hlchaptermodality{277}

\begin{chapteropening}
\textbf{By the end of this chapter:} \emph{I can describe the stages of an event: a step, a phase, the order of events, an instant, and what came at the end.}

Complete the last lesson of chapter 277: I can describe the stages of an event: a step, a phase, the order of events, an instant, and what came at the end.
\end{chapteropening}

\section[l'étape]{l'étape — the step, the leg; a stop on a journey}

\label{lesson:FR-C277-etape}



\begin{quote}
Before the new one: say the French for the era, then the French for the length.
\end{quote}

\subsection*{You'll want to know: l'étape}
\textbf{l'étape} — \textquotedblleft{}the step, the leg; a stop on a journey\textquotedblright{}.

\begin{quote}
La première étape du voyage était la plus longue.
\end{quote}

\begin{quote}
\emph{The first leg of the journey was the longest.}
\end{quote}

Also a stop on a trip: \emph{faire étape} — to stop over.

\begin{cousinweb}
\emph{étape} comes from Old French \emph{estaple} \textquotedblleft{}a market, a place to unload goods\textquotedblright{}, which was borrowed from Middle Dutch \emph{stapel} \textquotedblleft{}a pile, a depot\textquotedblright{} (EtymDB).
\end{cousinweb}

\begin{grammarlens}[title={What you've built}]
One more word for telling, explaining and describing at B1.
\end{grammarlens}

\subsection*{Your turn}
\begin{itemize}
\raggedright
  \item \emph{Say it:} \emph{l'étape}
  \item \emph{Say it:} \emph{l'étape}, once more
  \item \emph{Say it:} a sentence of your own with \emph{l'étape}
  \item \emph{Recall:} say the French for the era, then the French for the length, then say \emph{l'étape} again
\end{itemize}

\begin{tcolorbox}[breakable,colback=teal!4,colframe=teal!35!black,title={Before you move on}]
What does l'étape mean? (\textquotedblleft{}The step, the leg; a stop on a journey\textquotedblright{}.) Say it once more.
\end{tcolorbox}
\section[la phase]{la phase — the phase}

\label{lesson:FR-C277-phase}



\begin{quote}
Before the new one: say the French for the length, then the French for the step.
\end{quote}

\subsection*{You'll want to know: la phase}
\textbf{la phase} — \textquotedblleft{}the phase\textquotedblright{}.

\begin{quote}
Le projet entre dans sa dernière phase.
\end{quote}

\begin{quote}
\emph{The project is entering its last phase.}
\end{quote}

Also the phases of the moon: \emph{les phases de la lune}.

\begin{cousinweb}
\emph{phase} is a learned borrowing from Greek \emph{phásis} \textquotedblleft{}appearance\textquotedblright{}, said first of the moon (EtymDB lists the French word without a parent).
\end{cousinweb}

\begin{grammarlens}[title={What you've built}]
One more word for telling, explaining and describing at B1.
\end{grammarlens}

\subsection*{Your turn}
\begin{itemize}
\raggedright
  \item \emph{Say it:} \emph{la phase}
  \item \emph{Say it:} \emph{la phase}, once more
  \item \emph{Say it:} a sentence of your own with \emph{la phase}
  \item \emph{Recall:} say the French for the length, then the French for the step, then say \emph{la phase} again
\end{itemize}

\begin{tcolorbox}[breakable,colback=teal!4,colframe=teal!35!black,title={Before you move on}]
What does la phase mean? (\textquotedblleft{}The phase\textquotedblright{}.) Say it once more.
\end{tcolorbox}
\section[la chronologie]{la chronologie — the chronology, the order of events}

\label{lesson:FR-C277-chronologie}



\begin{quote}
Before the new one: say the French for the step, then the French for the phase.
\end{quote}

\subsection*{You'll want to know: la chronologie}
\textbf{la chronologie} — \textquotedblleft{}the chronology, the order of events\textquotedblright{}.

\begin{quote}
Le policier a refait la chronologie de la soirée.
\end{quote}

\begin{quote}
\emph{The police officer went back over the order of events that evening.}
\end{quote}

\emph{Par ordre chronologique} — in order of time.

\begin{cousinweb}
\emph{chronologie} is a learned word built from Greek \emph{khrónos} \textquotedblleft{}time\textquotedblright{} and \emph{-logía} \textquotedblleft{}account, study\textquotedblright{} (EtymDB lists the French word without a parent).
\end{cousinweb}

\begin{grammarlens}[title={What you've built}]
One more word for telling, explaining and describing at B1.
\end{grammarlens}

\subsection*{Your turn}
\begin{itemize}
\raggedright
  \item \emph{Say it:} \emph{la chronologie}
  \item \emph{Say it:} \emph{la chronologie}, once more
  \item \emph{Say it:} a sentence of your own with \emph{la chronologie}
  \item \emph{Recall:} say the French for the step, then the French for the phase, then say \emph{la chronologie} again
\end{itemize}

\begin{tcolorbox}[breakable,colback=teal!4,colframe=teal!35!black,title={Before you move on}]
What does la chronologie mean? (\textquotedblleft{}The chronology, the order of events\textquotedblright{}.) Say it once more.
\end{tcolorbox}
\section[l'instant]{l'instant — the instant, the moment}

\label{lesson:FR-C277-instant}



\begin{quote}
Before the new one: say the French for the phase, then the French for the chronology.
\end{quote}

\subsection*{You'll want to know: l'instant}
\textbf{l'instant} — \textquotedblleft{}the instant, the moment\textquotedblright{}.

\begin{quote}
Pendant un instant, personne n'a parlé.
\end{quote}

\begin{quote}
\emph{For a moment, nobody spoke.}
\end{quote}

\emph{Pour l'instant} — for now; \emph{à l'instant} — a moment ago.

\begin{cousinweb}
\emph{instant} is a learned borrowing from Latin \emph{īnstāns} \textquotedblleft{}present, pressing\textquotedblright{} (EtymDB), from \emph{īnstāre} \textquotedblleft{}to stand near, to press on\textquotedblright{} (Lewis \& Short).
\end{cousinweb}

\begin{grammarlens}[title={What you've built}]
One more word for telling, explaining and describing at B1.
\end{grammarlens}

\subsection*{Your turn}
\begin{itemize}
\raggedright
  \item \emph{Say it:} \emph{l'instant}
  \item \emph{Say it:} \emph{l'instant}, once more
  \item \emph{Say it:} a sentence of your own with \emph{l'instant}
  \item \emph{Recall:} say the French for the phase, then the French for the chronology, then say \emph{l'instant} again
\end{itemize}

\begin{tcolorbox}[breakable,colback=teal!4,colframe=teal!35!black,title={Before you move on}]
What does l'instant mean? (\textquotedblleft{}The instant, the moment\textquotedblright{}.) Say it once more.
\end{tcolorbox}
\section[à l'issue de]{à l'issue de — at the end of, following}

\label{lesson:FR-C277-a-l-issue-de}



\begin{quote}
Before the new one: say the French for the chronology, then the French for the instant.
\end{quote}

\subsection*{You'll want to know: à l'issue de}
\textbf{à l'issue de} — \textquotedblleft{}at the end of, following\textquotedblright{}.

\begin{quote}
À l'issue du match, les joueurs ont salué le public.
\end{quote}

\begin{quote}
\emph{At the end of the match, the players waved to the crowd.}
\end{quote}

Formal and frequent in news reports. \emph{L'issue} alone is the way out, or the outcome.

\begin{cousinweb}
\emph{issue} is the old past participle of Old French \emph{issir} \textquotedblleft{}to go out\textquotedblright{}, which is inherited from Latin \emph{exīre} \textquotedblleft{}to go out\textquotedblright{} (Lewis \& Short).
\end{cousinweb}

\begin{grammarlens}[title={What you've built}]
One more word for telling, explaining and describing at B1.
\end{grammarlens}

\subsection*{Your turn}
\begin{itemize}
\raggedright
  \item \emph{Say it:} \emph{à l'issue de}
  \item \emph{Say it:} \emph{à l'issue de}, once more
  \item \emph{Say it:} a sentence of your own with \emph{à l'issue de}
  \item \emph{Recall:} say the French for the chronology, then the French for the instant, then say \emph{à l'issue de} again
\end{itemize}

\begin{tcolorbox}[breakable,colback=teal!4,colframe=teal!35!black,title={Before you move on}]
What does à l'issue de mean? (\textquotedblleft{}At the end of, following\textquotedblright{}.) Say it once more.
\end{tcolorbox}
